\documentclass{article[11pt,letterpaper]}

% --- Packages ---
\usepackage[utf8]{inputenc}  % Input encoding
\usepackage[T1]{fontenc}  % Font encoding
\usepackage{lmodern}  % Latin Modern fonts
\usepackage{amsmath}  % Advanced mathematics
\usepackage{amsfonts}  % AMS fonts
\usepackage{amssymb}  % AMS symbols
\usepackage{graphicx}  % Include graphics
\usepackage{xcolor}  % Color support
\usepackage{booktabs}  % Professional tables
\usepackage{hyperref}  % Hyperlinks
\usepackage{caption}  % Caption customization
\usepackage{subcaption}  % Subfigures
\usepackage{tabularx}  % Extended tables
\usepackage{longtable}  % Multi-page tables
\usepackage{mathtools}  % Math enhancements
\usepackage[shortintertext]{nccmath}  % Medium-sized math
\usepackage{listings}  % Code listings
\usepackage[dvipsnames]{xcolor}  % More colors for listings

% --- Hyperref Configuration ---
\hypersetup{
    colorlinks=true,
    linkcolor=blue!60!black,
    citecolor=blue!60!black,
    urlcolor=blue!60!black,
    pdftitle={The Unified Synthonicon: A Predictive Grammar for Molecular, Supramolecular, and Temporal Architectures},
    pdfauthor={}
}

% --- Document Metadata ---
\title{The Unified Synthonicon: A Predictive Grammar for Molecular, Supramolecular, and Temporal Architectures}
\author{}
\date{\today}

% --- Custom Commands ---
\newcommand{\TODO}[1]{\textcolor{red}{[TODO: #1]}}


\begin{document}

% --- Title ---
\maketitle

% --- Abstract ---
% --- Table of Contents ---
\tableofcontents

\clearpage

% --- Main Content ---
\section{The Unified Synthonicon: A Predictive Grammar for Molecular, Supramolecular, and Temporal Architectures}
\label{sec:the-unified-synthonicon-a-predictive-grammar-for-m}

% Content for The Unified Synthonicon: A Predictive Grammar for Molecular, Supramolecular, and Temporal Architectures

\subsection{I. Formalizing the Synthon: An Information-Theoretic Definition}
\label{sec:i-formalizing-the-synthon-an-information-theoretic}

% Content for I. Formalizing the Synthon: An Information-Theoretic Definition

\subsection{II. The Eleven Primitives}
\label{sec:ii-the-eleven-primitives}

% Content for II. The Eleven Primitives

\subsubsection{Primitive notes}
\label{sec:primitive-notes}

% Content for Primitive notes

\subsubsection{\textbackslash{}$T\textbackslash{}_\textbackslash{}{\textbackslash{}square\textbackslash{}square\textbackslash{}}\textbackslash{}$ — Cage Topology: Formal Entry}
\label{sec:tsquaresquare--cage-topology-formal-entry}

% Content for $T_{\square\square}$ — Cage Topology: Formal Entry

\subsection{III. The Kinetic Primitive \textbackslash{}$K\textbackslash{}$: Separation of Thermodynamic and Kinetic Fidelity}
\label{sec:iii-the-kinetic-primitive-k-separation-of-thermody}

% Content for III. The Kinetic Primitive $K$: Separation of Thermodynamic and Kinetic Fidelity

\subsection{IV. Composition Axioms: The Grammar's Production Rules}
\label{sec:iv-composition-axioms-the-grammars-production-rule}

% Content for IV. Composition Axioms: The Grammar's Production Rules

\subsection{V. Theoretical Underpinnings: Constraint Propagation, Non-Equilibrium Thermodynamics, and the \textbackslash{}$\textbackslash{}xi\textbackslash{}_\textbackslash{}{CP\textbackslash{}}\textbackslash{}$ Metric}
\label{sec:v-theoretical-underpinnings-constraint-propagation}

% Content for V. Theoretical Underpinnings: Constraint Propagation, Non-Equilibrium Thermodynamics, and the $\xi_{CP}$ Metric

\subsubsection{Constraint propagation and dissipative structures}
\label{sec:constraint-propagation-and-dissipative-structures}

% Content for Constraint propagation and dissipative structures

\subsubsection{\textbackslash{}$\textbackslash{}eta\textbackslash{}_\textbackslash{}{CP\textbackslash{}}\textbackslash{}$ and \textbackslash{}$\textbackslash{}xi\textbackslash{}_\textbackslash{}{CP\textbackslash{}}\textbackslash{}$: constraint propagation efficiency and inefficiency index}
\label{sec:etacp-and-xicp-constraint-propagation-efficiency-a}

% Content for $\eta_{CP}$ and $\xi_{CP}$: constraint propagation efficiency and inefficiency index

\subsection{VI. The Criticality Condition and the G–D Phase Diagram}
\label{sec:vi-the-criticality-condition-and-the-gd-phase-diag}

% Content for VI. The Criticality Condition and the G–D Phase Diagram

\subsubsection{Quantitative degeneracy measurement (v2.2)}
\label{sec:quantitative-degeneracy-measurement-v22}

% Content for Quantitative degeneracy measurement (v2.2)

\subsection{VII. The Relational Substrate}
\label{sec:vii-the-relational-substrate}

% Content for VII. The Relational Substrate

\subsubsection{Directed Constraints, Not Mere Relations}
\label{sec:directed-constraints-not-mere-relations}

% Content for Directed Constraints, Not Mere Relations

\subsubsection{The Algebra Has No Unary Information Generators}
\label{sec:the-algebra-has-no-unary-information-generators}

% Content for The Algebra Has No Unary Information Generators

\subsubsection{Axioms as Relational Well-Formedness Constraints}
\label{sec:axioms-as-relational-well-formedness-constraints}

% Content for Axioms as Relational Well-Formedness Constraints

\subsubsection{Encodability Does Not Imply Isolation}
\label{sec:encodability-does-not-imply-isolation}

% Content for Encodability Does Not Imply Isolation

\subsubsection{Ordinal Information Is Sufficient}
\label{sec:ordinal-information-is-sufficient}

% Content for Ordinal Information Is Sufficient

\subsection{VIII. Occam Targets — Three Free Parameters Eliminated}
\label{sec:viii-occam-targets--three-free-parameters-eliminat}

% Content for VIII. Occam Targets — Three Free Parameters Eliminated

\subsubsection{1. \textbackslash{}$\textbackslash{}lambda\textbackslash{}$ Has No Free Value: It Is the Primitive Matching Fraction (P-20)}
\label{sec:1-lambda-has-no-free-value-it-is-the-primitive-mat}

% Content for 1. $\lambda$ Has No Free Value: It Is the Primitive Matching Fraction (P-20)

\subsubsection{2. F-Tier Boundaries Are Integer Boltzmann Discrimination Ratios (P-21)}
\label{sec:2-f-tier-boundaries-are-integer-boltzmann-discrimi}

% Content for 2. F-Tier Boundaries Are Integer Boltzmann Discrimination Ratios (P-21)

\subsubsection{3. \textbackslash{}$\textbackslash{}Omega\textbackslash{}$ Is a Derived Label, Not an Independent Primitive (P-22)}
\label{sec:3-omega-is-a-derived-label-not-an-independent-prim}

% Content for 3. $\Omega$ Is a Derived Label, Not an Independent Primitive (P-22)

\subsubsection{Summary of Eliminated Free Parameters}
\label{sec:summary-of-eliminated-free-parameters}

% Content for Summary of Eliminated Free Parameters

\subsection{IX. Tuple Algebra and Compositional Design}
\label{sec:ix-tuple-algebra-and-compositional-design}

% Content for IX. Tuple Algebra and Compositional Design

\subsubsection{VI.1 The Tuple Quasi-Metric}
\label{sec:vi1-the-tuple-quasi-metric}

% Content for VI.1 The Tuple Quasi-Metric

\subsubsection{VI.2 Lattice Operations: Meet and Join}
\label{sec:vi2-lattice-operations-meet-and-join}

% Content for VI.2 Lattice Operations: Meet and Join

\subsubsection{VI.3 Path Algebra: The Valid-Swap Directed Graph}
\label{sec:vi3-path-algebra-the-valid-swap-directed-graph}

% Content for VI.3 Path Algebra: The Valid-Swap Directed Graph

\subsubsection{VI.4 Tensor Product: Ensemble Prediction}
\label{sec:vi4-tensor-product-ensemble-prediction}

% Content for VI.4 Tensor Product: Ensemble Prediction

\subsubsection{VI.5 Natural Transformations: Cross-Domain Lifts}
\label{sec:vi5-natural-transformations-cross-domain-lifts}

% Content for VI.5 Natural Transformations: Cross-Domain Lifts

\subsubsection{VI.6 DesignPipeline: The Writer+Maybe Monad}
\label{sec:vi6-designpipeline-the-writermaybe-monad}

% Content for VI.6 DesignPipeline: The Writer+Maybe Monad

\subsubsection{VI.7 Compositionality and the Monad Stack}
\label{sec:vi7-compositionality-and-the-monad-stack}

% Content for VI.7 Compositionality and the Monad Stack

\subsection{X. Hybrid Systems and Programmable Matter}
\label{sec:x-hybrid-systems-and-programmable-matter}

% Content for X. Hybrid Systems and Programmable Matter

\subsubsection{X.1 Programmable Matter as Path Algebra}
\label{sec:x1-programmable-matter-as-path-algebra}

% Content for X.1 Programmable Matter as Path Algebra

\subsubsection{X.2 The Eleven Programmable Matter Encodings}
\label{sec:x2-the-eleven-programmable-matter-encodings}

% Content for X.2 The Eleven Programmable Matter Encodings

\subsubsection{X.3 The \textbackslash{}$F\textbackslash{}$–\textbackslash{}$K\textbackslash{}$ Programmability Quadrant}
\label{sec:x3-the-fk-programmability-quadrant}

% Content for X.3 The $F$–$K$ Programmability Quadrant

\subsubsection{X.4 \textbackslash{}$\textbackslash{}Phi\textbackslash{}_c\textbackslash{}$ as Global Programmability}
\label{sec:x4-phic-as-global-programmability}

% Content for X.4 $\Phi_c$ as Global Programmability

\subsubsection{X.5 The Programmability Lattice}
\label{sec:x5-the-programmability-lattice}

% Content for X.5 The Programmability Lattice

\subsubsection{X.6 Cross-Domain Analogies}
\label{sec:x6-cross-domain-analogies}

% Content for X.6 Cross-Domain Analogies

\subsection{XI. Experimental and Computational Validation Across Domains}
\label{sec:xi-experimental-and-computational-validation-acros}

% Content for XI. Experimental and Computational Validation Across Domains

\subsubsection{Molecular domain}
\label{sec:molecular-domain}

% Content for Molecular domain

\subsubsection{Supramolecular domain}
\label{sec:supramolecular-domain}

% Content for Supramolecular domain

\subsubsection{Temporal domain}
\label{sec:temporal-domain}

% Content for Temporal domain

\subsubsection{Network and bulk material domain}
\label{sec:network-and-bulk-material-domain}

% Content for Network and bulk material domain

\subsubsection{Quantum domain — Axiom 1 as a classical boundary detector}
\label{sec:quantum-domain--axiom-1-as-a-classical-boundary-de}

% Content for Quantum domain — Axiom 1 as a classical boundary detector

\subsubsection{Quantum domain — Topological matter (v0.4.0 catalog entries)}
\label{sec:quantum-domain--topological-matter-v040-catalog-en}

% Content for Quantum domain — Topological matter (v0.4.0 catalog entries)

\subsubsection{Quantum domain — Tuple-space phase diagram (v0.4.0)}
\label{sec:quantum-domain--tuple-space-phase-diagram-v040}

% Content for Quantum domain — Tuple-space phase diagram (v0.4.0)

\subsubsection{Planned validation}
\label{sec:planned-validation}

% Content for Planned validation

\paragraph{Transformation \textbackslash{}#8 — Literature-Grounded Partial \textbackslash{}$\textbackslash{}Phi\textbackslash{}_c\textbackslash{}$ Anchor (Groppi et al. 2020)}
\label{sec:transformation-8--literature-grounded-partial-phic}

% Content for Transformation #8 — Literature-Grounded Partial $\Phi_c$ Anchor (Groppi et al. 2020)

\subsubsection{Molecular domain — Design program suite (v0.4.1)}
\label{sec:molecular-domain--design-program-suite-v041}

% Content for Molecular domain — Design program suite (v0.4.1)

\section{synthomnicon/domains/molecular/\textbackslash{}_\textbackslash{}_init\textbackslash{}_\textbackslash{}_.py}
\label{sec:synthomnicondomainsmolecularinitpy}

% Content for synthomnicon/domains/molecular/__init__.py

\section{synthomnicon/\textbackslash{}_\textbackslash{}_init\textbackslash{}_\textbackslash{}_.py  (auto-registration chain)}
\label{sec:synthomniconinitpy--auto-registration-chain}

% Content for synthomnicon/__init__.py  (auto-registration chain)

\subsection{XII. Computational Validation Summary}
\label{sec:xii-computational-validation-summary}

% Content for XII. Computational Validation Summary

\subsection{XIII. Catalog and Audit Framework}
\label{sec:xiii-catalog-and-audit-framework}

% Content for XIII. Catalog and Audit Framework

\subsubsection{Catalog}
\label{sec:catalog}

% Content for Catalog

\subsubsection{Contamination and the 4-pass audit}
\label{sec:contamination-and-the-4-pass-audit}

% Content for Contamination and the 4-pass audit

\subsubsection{Analogy engine}
\label{sec:analogy-engine}

% Content for Analogy engine

\subsection{XIV. AI-Driven Design and Cross-Domain Similarity Search}
\label{sec:xiv-ai-driven-design-and-cross-domain-similarity-s}

% Content for XIV. AI-Driven Design and Cross-Domain Similarity Search

\subsection{XV. Future Extensions and Critical Considerations}
\label{sec:xv-future-extensions-and-critical-considerations}

% Content for XV. Future Extensions and Critical Considerations

\subsection{XVI. The SM/QG Disparity as a Primitive Mismatch}
\label{sec:xvi-the-smqg-disparity-as-a-primitive-mismatch}

% Content for XVI. The SM/QG Disparity as a Primitive Mismatch

\subsection{XVII. The Gravity Theory Spectrum: SM and QG Compatibility}
\label{sec:xvii-the-gravity-theory-spectrum-sm-and-qg-compati}

% Content for XVII. The Gravity Theory Spectrum: SM and QG Compatibility

\subsubsection{Encodings}
\label{sec:encodings}

% Content for Encodings

\subsubsection{Compatibility spectrum}
\label{sec:compatibility-spectrum}

% Content for Compatibility spectrum

\subsubsection{Key results}
\label{sec:key-results}

% Content for Key results

\subsection{XVIII. Three Stress Tests: Black Hole Entropy, Hilbert Space Factorization, and ER=EPR}
\label{sec:xviii-three-stress-tests-black-hole-entropy-hilber}

% Content for XVIII. Three Stress Tests: Black Hole Entropy, Hilbert Space Factorization, and ER=EPR

\subsubsection{XIX.1 Black Hole Entropy Scaling}
\label{sec:xix1-black-hole-entropy-scaling}

% Content for XIX.1 Black Hole Entropy Scaling

\subsubsection{XIX.2 Hilbert Space Factorization Failure}
\label{sec:xix2-hilbert-space-factorization-failure}

% Content for XIX.2 Hilbert Space Factorization Failure

\subsubsection{XIX.3 ER=EPR and R-Primitive Degeneracy at \textbackslash{}$G\textbackslash{}_\textbackslash{}{\textbackslash{}aleph\textbackslash{}}\textbackslash{}$}
\label{sec:xix3-erepr-and-r-primitive-degeneracy-at-galeph}

% Content for XIX.3 ER=EPR and R-Primitive Degeneracy at $G_{\aleph}$

\subsection{XIX. Development Roadmap}
\label{sec:xix-development-roadmap}

% Content for XIX. Development Roadmap

% --- Conclusion ---
\section{Conclusion}
\label{sec:conclusion}

% Conclusion content here

% --- Bibliography ---
% \bibliographystyle{plain}
% \bibliography{references}

\end{document}